% SOURCE_AUTHORITY: sga5_fr_workpass.tex, lines 7059--7086 (LF-normalized unit).
% SOURCE_SHA256: 791F4EFFC5E02832D5D77ED03518C8156D6F07E4C8238B03545DB93D883FBB28
\subsubsection*{2.5. Caracterización de los sistemas proyectivos que verifican la condición de Mittag-Leffler-Artin-Rees}
\begin{statement}{Proposición 2.5.1}
Las afirmaciones siguientes para un objeto $X$ de $P$ son equivalentes:
\begin{enumerate}[label=(\roman*)]
\item $X$ verifica (MLAR) (2.1.1).
\item Existe un sistema proyectivo estricto $Y$ y un AR-isomorfismo $u:Y\to X$ (2.4.5).
\item Existe un sistema proyectivo estricto $Y$ tal que $X$ e $Y$ sean isomorfos en la categoría $P_{AR}$.
\item $X$ verifica la condición de Mittag-Leffler y, denotando por $X'$ el sistema proyectivo de las imágenes universales de $X$, el sistema proyectivo $X/X'$ es AR-nulo.
\end{enumerate}
\end{statement}

\begin{prooftext}[Demostración]
Demostraremos la sucesión de implicaciones
\[
(i)\Rightarrow(iv)\Rightarrow(iii)\Rightarrow(ii)\Rightarrow(i).
\]
$(i)\Rightarrow(iv)$. Existe un entero $m$ tal que el morfismo canónico $w_{mX}$ de $X[m]$ en $X$ se factoriza a través de la inclusión canónica $i$ de $X'$ en $X$. De ello se deduce que $p_{AR}(i)$ es un epimorfismo escindido y, por tanto, puesto que $p_{AR}$ es exacto, que $p_{AR}(X/X')=0$. Se concluye mediante (2.4.4) (iii).

$(iv)\Rightarrow(iii)$. La condición (iv) significa, con la notación precedente, que $p_{AR}(i)$ es un epimorfismo. Pero, debido a la exactitud de $p_{AR}$, es además un monomorfismo y, por tanto, un isomorfismo. Se tomará como objeto $Y$ el sistema proyectivo $X'$.

$(iii)\Rightarrow(ii)$. El $P_{AR}$-isomorfismo de $Y$ en $X$ cuya existencia se supone es la imagen por $p_{AR}$ de un $P$-morfismo $Y[r]\to X$, donde $r$ designa un entero positivo o nulo. Se concluye gracias al hecho de que $Y[r]$ es estricto.

$(ii)\Rightarrow(i)$. Se puede suponer que $Y$ está contenido en $X$ y que $u$ es la inyección canónica. Sea $T$ el conúcleo de $u$. La hipótesis significa que $T$ es AR-nulo, por lo que existe un entero $r$ tal que, para todo $s$ mayor o igual que $r$, el morfismo canónico de $T[s]$ en $T$ sea nulo. Se concluye que, para $s\ge r$, la imagen de $X[s]$ en $X$ está contenida en $Y$; como, por otra parte, contiene evidentemente la imagen, igual a $Y$, de $Y[s]$ en $Y$, es igual a $Y$. Se concluye que $Y$ es la imagen de $X[s]$ para $s$ suficientemente grande, por lo que $X$ verifica (MLAR), y que $Y$ es el sistema proyectivo de las imágenes universales de $X$.
\end{prooftext}

\begin{statement}{Corolario 2.5.2}
Sea $u:Y\to X$ un AR-isomorfismo, siendo $Y$ estricto. En estas condiciones, $X$ verifica la condición (MLAR) y $u(Y)$ es el sistema proyectivo de las imágenes universales de $X$.
\end{statement}
